\documentclass[11pt]{article}
\usepackage[T1]{fontenc}
\usepackage{lmodern,amsmath,amssymb,amsthm,booktabs,microtype}
\usepackage[margin=1in]{geometry}
\usepackage[hidelinks]{hyperref}
\newtheorem{theorem}{Theorem}
\newtheorem{corollary}[theorem]{Corollary}
\newcommand{\ord}{\operatorname{ord}}
\newcommand{\Z}{\mathbb Z}
\newcommand{\Qzero}{q_0}
\title{An exact obstruction for the sparse places in the Lean proof of Dirichlet GRH}
\author{}
\date{}
\begin{document}
\maketitle
\vspace{-2em}
\begin{abstract}
The Lean construction of the distinguished infinite places in a cyclotomic
tower avoids triples formed by signed powers of a prime. We determine exactly
which primes can be admitted simultaneously. A single explicit Behrend-type
set works for every prime for which even a singleton can work. In one instance
we give 120 places and certify that the first obstructing prime is
$18\,216\,708\,648\,992\,767$. The argument is finite and independent of any
hypothesis about zeros of $L$-functions.
\end{abstract}

\section{The sparse-place problem}

Fix a multiple $q_0$ of $18480$, an integer $s\ge1$, and put
\[
h=3^{\lfloor\log_3 s\rfloor+4},\qquad
Q=q_0\cdot3h\cdot2^s,\qquad v=v_2(q_0),\qquad N=2^{s+v}.
\]
The archive identifies the place labels with $(\Z/Q\Z)^\times$. The desired
labels reduce to $-1$ modulo $q_0$. For a prime $p$, let $M_p=Q/p^{v_p(Q)}$.
A set $T$ is \emph{sparse at $p$} if there are no
$x,y,z\in T$, $0<i<j<h$, and signs $\epsilon,\delta\in\{1,-1\}$ with
\[
y/x\equiv \epsilon p^i\pmod{M_p},\qquad
z/x\equiv \delta p^j\pmod{M_p}.
\]
The entries can coincide, and the signs are independent. This is precisely
the local condition in \texttt{New.Parameters.IsSparse}. The upstream theorem
\texttt{Parameters.exists\_large\_isSparse} uses
$(2hH)^n\le2^s$ and $2h(L+1)<2^s$ to obtain at least
$H^n/(nH^2)$ sparse places for primes $p\le L$
\cite{definition,construction}.

Let $d_p$ denote the order of the class of $p$ in
$(\Z/M_p\Z)^\times/\{\pm1\}$. Equivalently, $d_p$ is the first positive
exponent $t$ for which $p^t\equiv\pm1\pmod{M_p}$.

\begin{theorem}\label{thm}
Let $H\ge2$, $n\ge1$, and $B=(h-1)(H-1)+1$. Assume $B^n\le2^s$.
There is a single set $T$ of labels $-1\pmod{q_0}$ with
\[
|T|\ \ge\ \frac{H^n}{nH^2}
\]
such that, for every prime $p$, the following three conditions are equivalent:
$T$ is sparse at $p$, some nonempty set is sparse at $p$, and $2d_p\ge h$.
The same equivalence holds simultaneously for every prescribed collection of
primes, with the numerical condition imposed separately on each prime.
\end{theorem}

\begin{proof}
A prime with $2d_p<h$ obstructs any nonempty set: use $x=y=z$,
$i=d_p$, $j=2d_p$. We construct a set for which all obstructions arise this way.

Take the digit cube $\{0,\ldots,H-1\}^n$, centered at
$c=((H-1)/2,\dots,(H-1)/2)$. All values of $\|u-c\|^2$ differ by integers,
since $a^2-(H-1)a$ is an integer for each digit $a$. Their range has length
at most $n(H-1)^2/4$. The most populated sphere therefore contains
at least $H^n/(1+\lfloor n(H-1)^2/4\rfloor)\ge H^n/(nH^2)$ points.
Encode each such $u$ by $e(u)=\sum_{k=0}^{n-1}u_k B^k$. Whenever
\[
j e(y)\equiv(j-i)e(x)+i e(z)\pmod{2^s},\qquad 0<i<j<h,
\]
the digit on either side is at most $(h-1)(H-1)=B-1$ at each position.
Both encoded sides are at most $B^n-1<2^s$. Equality of the encoded integers
then forces equality in every digit. The vectors satisfy
$y=(1-i/j)x+(i/j)z$. Strict convexity of the squared norm on the sphere
centered at $c$ forces $x=y=z$.

Put $g=1+3hq_0$ and $T=\{-g^{e(u)}\bmod Q: u\text{ on the chosen sphere}\}$.
The integer $g$ is prime to $Q$ and congruent to $1\pmod{q_0}$.
As $v_2(g-1)=v\ge4$,
$v_2(g^a-1)=v+v_2(a)$ for $a>0$; indeed an odd power preserves the
valuation, and each successive squaring adds one. Hence $g$ has order $2^s$
modulo $N$, and all these labels are distinct.

For odd $p$, the divisor $N$ of $M_p$ retains both congruences of any
forbidden triple. On $(\Z/N\Z)^\times$, send a unit $u$ to its unique signed
representative $r(u)\in\{u,-u\}$ that is $1\pmod4$. This is a group
homomorphism, it kills both signs, and it fixes $g$. Writing the exponents
as $e_1,e_2,e_3$, the congruences give
\[
 g^{e_2-e_1}=r(p)^i,\qquad g^{e_3-e_1}=r(p)^j\pmod N.
\]
Raise the first relation to $j$ and the second to $i$. The order of $g$
gives $j e_2\equiv(j-i)e_1+i e_3\pmod{2^s}$. The digit argument shows
$e_1=e_2=e_3$. Returning to the original full modulus $M_p$, both $i$ and
$j$ are multiples of $d_p$, so $2d_p\le j<h$.

For $p=2$, every ratio of two elements of $T$ is $1$ modulo $3h$, a divisor
of $M_2$. A forbidden triple would imply $4^i\equiv1\pmod{3h}$. The order of
$4$ modulo $3h$ is $h$, since $v_3(4^a-1)=1+v_3(a)$ for $a>0$.
This contradicts $i<h$. Together with the singleton observation, this also
gives $2d_2\ge h$.
\end{proof}

\begin{corollary}\label{cut}
The largest integer cutoff $L$ for which any nonempty set is sparse at every
prime $p\le L$ is
\[
\min\{p\text{ prime}:2d_p<h\}-1.
\]
The same $T$ from Theorem~\ref{thm} realizes this cutoff. In particular,
$h(L+1)\le2^{s+v+1}$ is sufficient.
\end{corollary}

\begin{proof}
A prime $p\equiv1\pmod Q$ has $d_p=1$; such primes exist by Dirichlet's
theorem. This proves existence of a first obstruction. For odd $p$, choose
$a=p$ or $-p$ so that $a\equiv1\pmod4$. Its order modulo $N$ is
$2^{\max(0,s+v-v_2(a-1))}$ by the same valuation formula.
This order divides $d_p$. Since $2^{v_2(a-1)}\le p+1$, we have
$d_p\ge 2^{s+v}/(p+1)$. For $p=2$ the proof of Theorem~\ref{thm} applies.
\end{proof}

\section{An exact finite instance}

Take $q_0=18480$, $s=44$ and $h=2187$. Then
\[
 Q=2^{48}3^9\cdot5\cdot7\cdot11
 =2\,133\,004\,707\,139\,399\,188\,480.
\]
The choice $s=44$ is compatible with the tower identity $2^s=m^{22}$, with
$m=4$. This example tests the finite sparse-place lemma alone.
With $n=3$, $H=12$ and $B=24047$ we have $B^3=13\,905\,375\,151\,823<2^{44}$.
Exactly 120 triples $(a,b,c)\in\{0,\ldots,11\}^3$ satisfy
\[
 (2a-11)^2+(2b-11)^2+(2c-11)^2=131.
\]
The unordered triples of odd squares are $(121,9,1)$, $(81,49,1)$ and
$(81,25,25)$. Signs and coordinate permutations give $48+48+24=120$
labels $-g^{a+Bb+B^2c}\bmod Q$.

The original sufficient numerical cutoff is $4\,021\,990\,406$, and
Corollary~\ref{cut}'s elementary bound gives $257\,407\,386\,108$. The
exact largest cutoff is $18\,216\,708\,648\,992\,766$. The following
certificate proves minimality of the next prime, $p_*$.

For an odd bad prime one must have $d_p\le1093$ and its projective order
modulo $2^{48}$ divides $d_p$. This dyadic order is a power of two at most
$1024$. Thus $p\equiv\pm1\pmod{2^{38}}$. Of the $132\,543$ integers
at least $2$ in these two classes through $p_*$, exactly $55\,083$ are prime
to $Q$. The primes dividing $Q$ other than $2$ all fail the two-class test.
For each retained integer $u$, calculate its dyadic order $d_2(u)$ as in
Corollary~\ref{cut}. Test the multiples $k$ of $d_2(u)$ through $1093$ for
$u^k\equiv\pm1\pmod Q$. Exactly seven integers pass.

\begin{center}
\begin{tabular}{rrl}
\toprule
Integer & Least exponent & Certificate\\
\midrule
$1\,020\,346\,790\,576\,129$ &1080& divisible by $157$\\
$2\,185\,829\,116\,018\,687$ &320& divisible by $101999$\\
$3\,870\,280\,929\,771\,521$ &972& divisible by $1663$\\
$6\,412\,351\,813\,189\,631$ &864& divisible by $17$\\
$9\,399\,724\,905\,857\,023$ &768& divisible by $23$\\
$17\,099\,604\,835\,172\,353$ &324& divisible by $13$\\
$18\,216\,708\,648\,992\,767$ &864& Lucas certificate, base $3$\\
\bottomrule
\end{tabular}
\end{center}

For the last integer $p_*$, the factorization of $p_*-1$ is
$2\cdot3^6\cdot11\cdot74413\cdot15264089$. Trial division proves primality
of the factors. The accompanying program checks that $3^{p_*-1}=1\pmod{p_*}$
and $\gcd(3^{(p_*-1)/r}-1,p_*)=1$ for each prime factor $r$.
These congruences give multiplicative order $p_*-1$ modulo every prime
dividing $p_*$, proving primality. At this prime $i=864$ and $j=1728$
obstruct even a singleton.

\section{Scope}

The construction of convex-position digit vectors belongs to the Behrend
tradition. Shapira treats simultaneous systems \cite{shapira}. A previous
ProveIt report \cite{previous} uses a digit sphere to recover the exact spectra
of several rational progression shapes in a fixed prime field. The present
problem has two independent signs and varies the prime-to-$p$ modulus with
the multiplier. The exact obstruction above depends on keeping the
sign quotient in the dyadic projection and then returning to the full
prime-to-$p$ modulus after the three exponents coincide. It solves the
optimal prime-cutoff question for the repository's specific sparse-place
condition and gives a stronger quantitative input to that proof. The largest
possible cardinality at a fixed cutoff is a separate extremal problem.
The global Dirichlet GRH argument still has all its other quantitative
inputs; no new zero-free theorem follows here.

\begin{thebibliography}{9}
\bibitem{definition} \texttt{lsernau/grh-lean},
\href{https://github.com/lsernau/grh-lean/blob/72fead102f663d5c7a407964a13dc0cf9d977207/grh/src/New/Reduction/Data/AdmissibleData.lean}{\texttt{New/Reduction/Data/AdmissibleData.lean}},
definition \texttt{New.Parameters.IsSparse}. Readable edition.
\bibitem{construction} Same repository and commit,
\href{https://github.com/lsernau/grh-lean/blob/72fead102f663d5c7a407964a13dc0cf9d977207/grh/src/New/Reduction/Data/Existence.lean}{\texttt{New/Reduction/Data/Existence.lean}},
\texttt{exists\_three\_term\_free} and \texttt{Parameters.exists\_large\_isSparse},
lines 80 and 167.
\bibitem{previous} ProveIt, \href{https://github.com/VladimirReshetnikov/ProveIt/blob/d31039a48d963692e5a5176f47b6603fc9e9d9cf/docs/incoming/dirichlet-grh-research.zip}{\texttt{dirichlet-grh-research.zip}}, \texttt{progressions/syndetic.tex}, Theorem 1.
\bibitem{shapira} A.~Shapira, \emph{Behrend-type constructions for sets of
linear equations}, Acta Arith. 122 (2006), 17--33.
\href{https://doi.org/10.4064/aa122-1-2}{doi:10.4064/aa122-1-2}.
\end{thebibliography}
\end{document}
